% FaithLM for Vietnamese — ACL short paper (4 pages + unlimited references)
%
% Build:  cd paper && bash build.sh          (or: latexmk -pdf main.tex)
% Numbers: python scripts/export_paper_tables.py    -> regenerates generated/*.tex
%
% Anything marked TODO needs a human. Everything inside generated/ is produced
% from the result JSON and should never be edited by hand.

\documentclass[11pt]{article}

% "review" = anonymous submission; switch to "final" for the version you submit
% to the course, so that author names are visible.
\usepackage[review]{acl}

\usepackage{times}
\usepackage{latexsym}

% T5 encoding renders Vietnamese diacritics correctly. Keep both lines: T1 must
% come first so that Latin text and the bibliography still hyphenate properly.
\usepackage[T1,T5]{fontenc}
\usepackage[utf8]{inputenc}

\usepackage{microtype}
\usepackage{inconsolata}
\usepackage{graphicx}
\usepackage{booktabs}
\usepackage{amsmath}

\title{TODO: Title --- e.g.\ ``How Faithful Are LLM Explanations in Vietnamese?
       Reproducing and Correcting FaithLM on XCOPA-vi''}

% For the course submission, switch \usepackage[review]{acl} to [final] and fill
% in the real names, student IDs and emails below.
\author{
  First Author \\
  MSHV: TODO \\
  \texttt{TODO@example.com} \\\And
  Second Author \\
  MSHV: TODO \\
  \texttt{TODO@example.com} \\\And
  Third Author \\
  MSHV: TODO \\
  \texttt{TODO@example.com} \\\And
  Fourth Author \\
  MSHV: TODO \\
  \texttt{TODO@example.com} \\
}

\begin{document}
\maketitle

\begin{abstract}
% ~150 words. Write this LAST, once the numbers exist.
% Suggested skeleton — replace with real findings, do not leave hedged text:
%   1 sentence: what faithfulness means and why it matters for Vietnamese.
%   1 sentence: what we did (reproduce FaithLM on XCOPA-vi with Qwen3-4B + DeepSeek).
%   2 sentences: the two defects we found in the released implementation.
%   2 sentences: what our corrected measurements show.
%   1 sentence: what this implies for evaluating explanations in non-English settings.
TODO: Abstract.
\end{abstract}

\section{Introduction}
\label{sec:intro}

% ~0.75 column. Cover, in order:
%   - Explanations from LLMs are increasingly used as evidence of reasoning, but
%     a fluent explanation need not be the reason the model answered as it did.
%   - FaithLM \citep{faithlm} tests this by perturbation: negate the explanation,
%     feed it back, and see whether the prediction moves.
%   - Nearly all of this evidence is English. Vietnamese is under-studied, and
%     it is not obvious that the method transfers.
%   - State the contributions explicitly (the list below).
TODO: Introduction.

Our contributions are:

\begin{itemize}
  \item We reproduce FaithLM on Vietnamese XCOPA with an open 4B predictor and
        an API explainer, and release a configuration-driven implementation.
  \item We identify two defects in the released implementation that affect the
        reported metric, and quantify the effect of correcting each.
  \item We compare against non-LLM counterfactual baselines, isolating how much
        of the measured faithfulness requires an LLM at all.
\end{itemize}

\section{Related Work}
\label{sec:related}

% ~0.5 column, three short paragraphs. Cite real work only — do not invent
% citations. Add entries to custom.bib as you go.
%   - Faithfulness of explanations: distinguish plausibility from faithfulness.
%   - Perturbation-based evaluation, and prompt optimisation by LLM (OPRO).
%   - Cross-lingual evaluation and XCOPA; note the scarcity of Vietnamese work.
TODO: Related work.

\section{Dataset}
\label{sec:dataset}

% ~0.4 column. State plainly:
%   - XCOPA-vi: source, size, split used, and that it is a human translation of
%     COPA rather than natively authored Vietnamese (this matters for Limitations).
%   - The 200-question subsample, drawn with a fixed seed over the full split.
%   - The English COPA control, used to separate "the method is weak" from
%     "the method is weak in Vietnamese".
TODO: Dataset description.

% Generated by scripts/export_paper_tables.py -- do not edit by hand.
% Rerun that script to refresh these numbers.

\begin{table}[t]
\centering
\small
\begin{tabular}{lrr}
\toprule
\textbf{Dataset} & \textbf{Split} & \textbf{Size} \\
\midrule
XCOPA-vi & test & 500 \\
XCOPA-vi (evaluated) & test & -- \\
COPA-en (control) & train & 1{,}000 \\
\bottomrule
\end{tabular}
\caption{Datasets. The evaluated subset is drawn from the full split with a
fixed seed, so every run covers the same questions. TODO: update caption.}
\label{tab:data}
\end{table}


\section{Method}
\label{sec:method}

\subsection{Faithfulness by counterfactual perturbation}

Let $q$ be a question, $\hat{y}$ the predictor's answer, and $e$ an explanation
of $\hat{y}$ produced by a separate explainer model. Let $\bar{e}$ be a
counterfactual explanation with the opposite meaning. Writing $s(\cdot)$ for the
probability the predictor assigns to the gold answer under a given hint, the
faithfulness score is the gap between the two arms:

\begin{equation}
\label{eq:faith}
  F = \left| s(\text{arm}_{\text{true}}) - s(\bar{e}) \right|
\end{equation}

% Explain in prose: a faithful explanation is one whose reversal moves the
% prediction, so a large gap means the explanation tracked the actual decision.

\subsection{Two definitions of the true arm}
\label{sec:two-metrics}

% This subsection is the paper's main analytical contribution. Say clearly:
%   - PAPER: arm_true is the UNHINTED baseline. The true explanation is generated
%     but never shown to the predictor, so F conflates the effect of adding any
%     hint with the effect of reversing meaning.
%   - SYMMETRIC: arm_true is the true explanation itself. Both arms are hinted,
%     so F isolates meaning reversal.
%   - We report both. Do not claim one is "correct" without the numbers backing it.
TODO: Explain the two metric variants.

\subsection{Continuous scoring}
\label{sec:scoring}

% Explain why this is necessary, not merely nicer:
%   - Scoring a single question by string match yields s in {0,1}, so F in
%     Eq. (1) is almost always exactly 0.
%   - The optimiser is handed a near-constant objective and cannot improve.
%   - We score by length-normalised log-probability over the answer choices,
%     giving continuous s from one forward pass and no extra compute.
TODO: Explain log-probability scoring.

\subsection{Iterative refinement}

% Brief. Explainer acts as optimiser over past (explanation, score) pairs.
% Local refines the explanation per question; global searches one shared
% instruction. Keep this short — it is reproduced from the original paper.
TODO: Describe the optimisation loop.

\section{Experimental Setup}
\label{sec:setup}

% Generated by scripts/export_paper_tables.py -- do not edit by hand.
% Rerun that script to refresh these numbers.

\begin{table}[t]
\centering
\small
\begin{tabular}{ll}
\toprule
\textbf{Component} & \textbf{Setting} \\
\midrule
Predictor & \texttt{--} \\
Explainer & \texttt{--} \\
Scorer & -- \\
Optimisation iterations & -- \\
Question sampling & -- (seed --) \\
\bottomrule
\end{tabular}
\caption{Experimental configuration, read back from the stored run config.
TODO: update caption.}
\label{tab:setup}
\end{table}


% Add anything the generated table cannot know: hardware used, wall-clock time,
% total API cost, and the exact commit hash the numbers came from.
TODO: Hardware, runtime and reproducibility details.

\section{Results and Discussion}
\label{sec:results}

% Generated by scripts/export_paper_tables.py -- do not edit by hand.
% Rerun that script to refresh these numbers.

\begin{table}[t]
\centering
\small
\begin{tabular}{lrrrrr}
\toprule
\textbf{Setting} & \textbf{N} & \textbf{Acc.} & \textbf{Faith.} & \textbf{SD} & \textbf{Zero} \\
\midrule
FaithLM (paper metric) & -- & -- & -- & -- & -- \\
FaithLM (symmetric) & -- & -- & -- & -- & -- \\
COPA-en (control) & -- & -- & -- & -- & -- \\
\bottomrule
\end{tabular}
\caption{Task accuracy and faithfulness on Vietnamese XCOPA. \textbf{Faith.} is
the mean best faithfulness score per question; \textbf{Zero} is the fraction of
questions scoring exactly zero. TODO: replace this caption once the numbers are in.}
\label{tab:main}
\end{table}


% Interpret, do not merely restate the table. Answer these questions in prose:
%   - Does FaithLM beat the non-LLM baselines? By how much?
%   - How far apart are the paper and symmetric metrics, and what does the gap
%     mean? If they agree, that is also a finding worth stating.
%   - Vietnamese vs English: is the gap in task accuracy, in faithfulness, or both?
%   - Does the optimisation loop actually improve scores over iterations?
TODO: Discussion of results.

% Generated by scripts/export_paper_tables.py -- do not edit by hand.
% Rerun that script to refresh these numbers.

\begin{table}[t]
\centering
\small
\begin{tabular}{lrrr}
\toprule
\textbf{Counterfactual source} & \textbf{Faith.} & \textbf{SD} & \textbf{N} \\
\midrule
FaithLM (paper metric) & -- & -- & -- \\
Negation baseline & -- & -- & -- \\
Identity baseline & -- & -- & -- \\
Shuffle baseline & -- & -- & -- \\
\bottomrule
\end{tabular}
\caption{Replacing the LLM counterfactual generator with rule-based baselines.
The identity baseline returns the explanation unchanged, so its score is the
floor the metric assigns to a non-reversal. TODO: update caption.}
\label{tab:baselines}
\end{table}


\section{Error Analysis}
\label{sec:error}

% Generated by scripts/export_paper_tables.py -- do not edit by hand.
% Rerun that script to refresh these numbers.

\begin{table}[t]
\centering
\small
\begin{tabular}{lrr}
\toprule
\textbf{Category} & \textbf{Count} & \textbf{Share} \\
\midrule
Correct, faithful & -- & -- \\
Correct, unfaithful & -- & -- \\
Wrong, faithful & -- & -- \\
Wrong, unfaithful & -- & -- \\
\bottomrule
\end{tabular}
\caption{Predictions cross-tabulated against explanation faithfulness
(threshold $0.1$). ``Correct, unfaithful'' is the case of interest: the answer is
right, but the explanation does not track the decision. TODO: update caption.}
\label{tab:errors}
\end{table}


% The interesting cell is "correct prediction, unfaithful explanation" — the model
% got the right answer for reasons the explanation does not capture. Quote one or
% two real examples with the Vietnamese text, and say what went wrong.
% Also report any systematic failures: refusals, unparseable outputs, degenerate
% counterfactuals (e.g. the explainer merely prepending a negation).
TODO: Error analysis with concrete examples.

\section{Conclusion}
\label{sec:conclusion}

% ~0.25 column. What was measured, what it means for Vietnamese specifically,
% and one concrete next step. No new claims here.
TODO: Conclusion.

% Limitations is mandatory at *ACL venues and does not count toward the page
% limit. It is a section*, so it is unnumbered.
\section*{Limitations}

% Be specific and honest — vague limitations sections read as boilerplate.
% Real limitations of this work, to expand into prose:
%   - XCOPA-vi is translated from English COPA, so it may carry translationese
%     and does not represent natively authored Vietnamese causal language.
%   - We evaluate N questions with M iterations, well below the original scale,
%     bounded by compute and API budget. State N and M explicitly.
%   - Log-probability scoring requires a local model with choice options; it does
%     not apply to API predictors or to open-ended tasks.
%   - The explainer is a closed API model whose version may change, which limits
%     exact reproducibility.
%   - A single predictor at one size; results may not hold at other scales.
%   - Single seed for the sampling draw; we do not report variance across seeds.
TODO: Limitations.

\section*{Ethical Considerations}

% Short but not empty. Cover:
%   - Faithfulness scores measure sensitivity to perturbation, not truthfulness.
%     A high score does not license presenting explanations to end users as the
%     model's actual reasoning.
%   - Data is public and contains no personal information.
%   - Note the API provider used and that prompts leave the local machine.
TODO: Ethical considerations.

\section*{Acknowledgments}

TODO: Acknowledgments (remove for the anonymous version).

\bibliography{custom}

\appendix

\section{Prompt Templates}
\label{sec:appendix-prompts}

% Appendices do not count toward the page limit. Paste the exact prompts from
% faithlm/prompts.py so the work is reproducible.
TODO: Prompt templates.

\section{Additional Results}
\label{sec:appendix-results}

% Generated by scripts/export_paper_tables.py -- do not edit by hand.
% Rerun that script to refresh these numbers.

\begin{table}[t]
\centering
\small
\begin{tabular}{lrrrr}
\toprule
\textbf{Variant} & \textbf{N} & \textbf{Best} & \textbf{Iter.\ 1} & \textbf{Iter.\ last} \\
\midrule
FaithLM (paper metric) & -- & -- & -- & -- \\
FaithLM (symmetric) & -- & -- & -- & -- \\
Negation baseline & -- & -- & -- & -- \\
Identity baseline & -- & -- & -- & -- \\
Shuffle baseline & -- & -- & -- & -- \\
COPA-en (control) & -- & -- & -- & -- \\
\bottomrule
\end{tabular}
\caption{All variants, including the change from the first to the last
optimisation iteration. TODO: update caption.}
\label{tab:per-variant}
\end{table}


\end{document}
